% 第7章 自然语言处理与大模型
\section{自然语言处理与大模型：从情绪到交易智能体}\label{sec:nlp}

文本是金融市场里信噪比最微妙的一类数据：单条新闻的噪声极大，但它承载的信息（基本面变化、预期修正、情绪）是价量数据里没有的。这一族的演进，恰好可以按“信息抽取的粒度”分成三代——词典、编码器、大模型，每一代都扩展了可用的信息，也带来了新的、更隐蔽的泄漏风险。

\subsection{第一代：情绪与词典}

Tetlock（2007）\cite{tetlock2007} 用报纸专栏的用词来度量投资者情绪，发现负面措辞与随后的价格下行相关，是“文本 $\to$ 交易信号”的早期实证。Baker 与 Wurgler（2006）\cite{bakerwurgler2006} 构造了综合情绪指数并检验其与横截面收益的关系。这一代的技术核心是词典法：Loughran 与 McDonald（2011）\cite{loughranmcdonald2011} 指出通用情感词典在金融文本上会严重误判（“liability”“tax”“cost”在金融语境下中性，在通用词典里可能是负面），并构建了金融专用情感词典，成为此后几乎所有金融文本研究的标准工具。

这一代的问题很清楚：\textbf{词典无法处理否定、语境与隐含表达}，它能抓到的是“词的极性”，而不是“句子的意思”。

\subsection{第二代：预训练编码器}

FinBERT 类模型把金融文本理解推进到上下文级别。Araci（2019）\cite{araci2019finbert}、Yang 等（2020）\cite{yang2020finbert} 与 Huang、Wang、Yang（2023）\cite{huang2022finbert} 分别在金融语料上继续预训练 BERT，在情感分类与信息抽取任务上超过通用模型与词典法。这一代的贡献是双重的：它提升了文本信号的质量，也确立了“先在金融语料上预训练、再在下游任务上微调”的范式。但它仍受限于标注数据、任务特定的输出格式（分类标签），且\textbf{只能处理被专门训练过的任务}。

\subsection{第三代：大模型与交易智能体}

大语言模型把这一族推向两个新方向。

\textbf{（一）LLM 直接抽取收益信号}。Lopez-Lira 与 Tang（2023）\cite{lopezlira2023} 用 ChatGPT 读公司新闻标题，把输出映射成交易信号并回测，报告了显著的多空收益。Ke、Kelly 与 Xiu（2019）\cite{kebeixiu2019} 更早地用文本数据预测收益，是把文本转成收益信号的代表性实证。这类工作的吸引力在于\textbf{无需标注数据、可零样本迁移}，但它的最大隐患也在这里：\textbf{如果用于回测的新闻落在模型的预训练语料覆盖期内，模型“知道”结果，回测收益就是数据泄漏的产物。}一个设计严谨的 LLM 回测，必须要么使用模型训练截止日之后的样本，要么严格证明测试文本不在训练语料中。这一点在 2023--2025 年的大量 LLM 交易论文中被系统性忽视，是本文审计（第 \ref{sec:audit} 章）里最严重的一类污染。

金融专用大模型（BloombergGPT\cite{wu2023bloomberggpt}、FinGPT\cite{yang2023fingpt} 等）试图通过领域预训练提升金融任务表现，其价值在信息抽取类任务上较明确，在“直接预测收益”上的证据则弱得多。

\textbf{（二）LLM 交易智能体}。FinMem\cite{yu2023finmem}、FinAgent、TradingAgents、StockAgent 等工作把 LLM 当作一个能记忆、能推理、能调用工具并做多步决策的智能体，让它在“看新闻--查指标--做决策--复盘”的循环里自主交易。这类工作是当前研究最热的区域，故事也最动听。近期也出现了一批试图把评测标准化的基准与长周期检验工作（StockBench\cite{stockbench2025}、InvestorBench\cite{investorbench2024}、长周期对照\cite{llmlongrun2025}），方向对，但整体仍处早期。本文对其持保留态度，理由有三：其一，\textbf{评测协议普遍薄弱}，很多工作只在小样本、短区间上报告，且缺少与强基线的公平对照；其二，\textbf{数据泄漏难以排除}，LLM 的预训练语料若覆盖回测期，智能体的“先知”能力无法与“推理”能力分离；其三，\textbf{成本与可复现性问题突出}，LLM 的推理成本高、API 版本漂移导致同一代码在不同时间给出不同结果，直接破坏可复现性。

\subsection{落地要素}

\begin{itemize}
  \item \textbf{严格的时间戳对齐}。所有文本必须带“市场首次可见时间”，并按该时间进入策略；用发布时间之前的信息是前瞻，用之后但含糊处理同样有问题。
  \item \textbf{训练语料截止日的红线}。任何用 LLM 做回测的工作，必须报告模型的知识截止日与样本期的关系，否则结论不可信。
  \item \textbf{幻觉与可追溯性}。LLM 生成的“事实”需要回到原文核验；用于交易的抽取结果应可溯源到具体文本片段。
  \item \textbf{成本核算}。逐条新闻调用 LLM 的成本远高于一次向量化，必须计入策略净收益。
\end{itemize}

\subsection{复现状态与失效边界}

\textbf{具有讽刺意味的是，这一族是六族里可复现性最差的。}词典法（Loughran--McDonald）与开源编码器（FinBERT 系列）的复现状态较好，因为词典与权重是固定的；而所有依赖商业 LLM 的工作，都面临“模型升级导致结果变化”的系统性问题——用 GPT-3.5 得到的结论，在 GPT-4 上未必复现，却没人愿意公开承认这一点。因此本文把“LLM 智能体交易”的整体证据强度定为\textbf{低}，并建议任何以此为据的落地尝试，必须自建前向（forward）样本——只相信模型截止日之后的数据。

失效边界除通用三条（信息耗尽、非平稳、成本）外，这一族还多一条{\textbf{泄漏}：LLM 知道答案。}本章要记住的一句话是：\textbf{大模型让“读懂文本”这件事的门槛大幅下降，但它同时让“回测里有多少是推理、有多少是背诵”变得极难分辨——在证明自己没有泄漏之前，任何 LLM 交易回测都只能算一个假设。}
